% SCOPE: 4.1 setup/protocol; 4.2 RQ1/T1 main accuracy (FILLED, preliminary
% single-seed gs300, robust math_verify re-grade) + degeneration-mitigation
% finding; 4.3 RQ2/T2 mechanism decomposition (Gandhi-arm + shift-only arm)
% = PENDING with pre-registered decision rules; 4.4 RQ3/T3 difficulty
% stratification (FILLED); 4.5 RQ4/T4 calibration (FILLED, meta-only).
% All FILLED cells are a SINGLE-SEED preliminary gs300 evaluation; the RQ2
% decomposition and multi-seed replication remain pending.
\section{Experiments}
\label{sec:experiments}

Our experiments answer four questions, each with a dedicated table:
does the \pmishift{} package improve accuracy over a confound-free
matched baseline (RQ1, Table~\ref{tab:t1}); where does the effect come
from (RQ2, Table~\ref{tab:t2}); how does it vary with difficulty and
distribution shift (RQ3, Table~\ref{tab:t3}); and what happens to
calibration (RQ4, Table~\ref{tab:t4}). The accuracy, stratification, and
calibration tables (T1, T3, T4) report a \emph{single-seed, preliminary}
gs300 evaluation under a robust \texttt{math\_verify} re-grade; we mark
every such number \emph{preliminary} and defer multi-seed replication.
The mechanism decomposition (RQ2/T2) is still \emph{pending}: its cells
are placeholders (\todo{}) governed by pre-registered decision rules
(Section~\ref{sec:exp:rq2}).

\subsection{Setup and evaluation protocol}
\label{sec:exp:setup}

\paragraph{Models and training.} All arms start from Qwen3-8B. The
treatment arm follows the pipeline of Section~\ref{sec:method}: SFT-1
primes \metaopen{}\dots\metaclose{} blocks inside the reasoning trace
(\texttt{v8\_meta\_inside\_strict} data), SFT-2 adds functional
redirect/verify behaviors (\texttt{rv\_functional}), and RL is a
GRPO-family objective (\texttt{TRIOBJ\_DCPO\_V4})~\citep{shao2024deepseekmath}
whose launched configuration is a \emph{five-head package}, not the
\pmishift{} reward in isolation: the correctness reward
$R_{\mathrm{corr}} \in \{-1,+1\}$ (\texttt{math\_verify}) with weight
$w_{\mathrm{corr}} = 1.0$; the \pmishift{} meta-region reward
$R_{\mathrm{shift}}$ ($w_{\mathrm{meta}} = 0.8$, warmed up over 80 steps);
and the protocol-maintenance heads and costs of
Section~\ref{sec:method:pmishift} ($w_{\mathrm{cal}} = 0.3$,
$w_{\mathrm{format}} = 0.35$, $w_{\mathrm{emit}} = 0.1$, length cost
$0.08$, truncation penalty $0.3$). Because all five heads move together,
the T1 treatment$-$base contrast estimates the effect of the
\emph{whole package} versus vanilla GRPO, and isolating the \pmishift{}
reward is exactly the job of RQ2 (Section~\ref{sec:exp:rq2}). The control
arm is the matched-base twin of Section~\ref{sec:method:twin}: identical
data with meta blocks stripped, identical non-meta hyperparameters
(byte-audited config diff), and vanilla GRPO with correctness-only reward.
The base arm emits \emph{no} meta blocks and \emph{no} confidence token
(0\% meta emission), a structural asymmetry we return to in RQ4. Both arms
are trained with a 4{,}096-token response cap and evaluated at both 4k and
16k (below).

\paragraph{Evaluation protocol.} We evaluate on a held-out suite of
1{,}030 problems: 500 GSM8K, 500 MATH500, and 30 AIME. Answers are graded
with a robust \texttt{math\_verify} re-grade (correcting the earlier
brittle grader that spuriously deflated correct answers); generation uses
temperature $0.7$ at both a 4k and a 16k max-token budget. We report avg@8
accuracy (avg@16 on AIME to reduce variance on the small set). Both arms
are evaluated \emph{in the same job with identical seeds}. Significance is
assessed by \emph{problem-paired bootstrap} over the held-out set (and, on
the $n{=}30$ AIME set, additionally by McNemar's test). The current gs300
results are a \emph{single training seed} and are therefore preliminary;
multi-seed mean$\pm$std replication is pending.

\paragraph{Stratified reporting.} All meta-vs-base comparisons are also
reported \emph{difficulty-stratified} (Section~\ref{sec:exp:rq3}). The
motivation is empirical: the accuracy payoff of the meta package is
strongly difficulty-dependent (saturated on easy quartiles, largest on
the hardest), whereas meta-block \emph{emission} is essentially
\emph{non-selective}---the treatment arm emits a block on 86--94\% of
responses across \emph{all} difficulty quartiles
(Section~\ref{sec:exp:rq3}). This dissociation---uniform emission but
difficulty-concentrated payoff---means the model over-emits on easy
problems (an efficiency cost we flag as a limitation), and it is why we
rely on the between-arm matched-base contrast, which compares the two arms
on the full identical evaluation set rather than conditioning on the
endogenous decision to emit.

\subsection{RQ1: Does the \pmishift{} package improve held-out accuracy?}
\label{sec:exp:rq1}

\paragraph{Question and design.} RQ1 asks whether the \pmishift{} package
improves end-task accuracy over an otherwise identical GRPO run.
Table~\ref{tab:t1} compares the matched base twin against the \pmishift{}
arm at the final checkpoint (gs300). Because the twin differs from the
treatment only in the meta package of Section~\ref{sec:method:twin}, any
gap is attributable to that package \emph{as a whole}; apportioning it to
the \pmishift{} reward specifically is RQ2.

\paragraph{Result.} The package improves accuracy on every domain, with
the gain growing sharply with domain difficulty. On GSM8K, \pmishift{}
scores $93.9$ vs.\ the base's $89.9$ ($+4.0$\,pp, problem-paired bootstrap
$p<.001$); on MATH500, $81.5$ vs.\ $62.7$ ($+18.8$\,pp, $p<.001$); on
AIME (avg@16, 16k budget), $18.5$ vs.\ $4.8$ ($+13.8$\,pp, bootstrap
$p<.001$; McNemar $p\in[.03,.06]$ on $n{=}30$). The sample-weighted
overall accuracy is $85.7$ vs.\ $74.2$ ($+11.5$\,pp). These are
single-seed preliminary numbers; the effect size and its concentration on
hard problems (RQ3) are the two robust qualitative signals.

\begin{table}[t]
\caption{\textbf{T1: main results} at gs300 on the held-out
1{,}030-problem suite (avg@8; AIME avg@16 at the 16k budget), robust
\texttt{math\_verify} re-grade. \emph{Preliminary, single seed.} The
\pmishift{} arm is the five-head package (correctness $+$ \pmishift{}
$+$ calibration/format/emission heads $+$ costs); the treatment$-$base
gap is the package effect vs.\ vanilla GRPO, decomposed in RQ2. Overall
is the 1{,}030-sample-weighted mean. $p$-values: problem-paired bootstrap;
$^{\dagger}$AIME additionally McNemar $p\in[.03,.06]$, $n{=}30$.}
\label{tab:t1}
\begin{center}
\small
\begin{tabular}{lcccc}
\toprule
Arm & GSM8K & MATH500 & AIME & Overall \\
\midrule
Qwen3-8B (raw)              & \todo{} & \todo{} & \todo{} & \todo{} \\
SFT-only (SFT-1+2)          & \todo{} & \todo{} & \todo{} & \todo{} \\
Matched base (vanilla GRPO) & $89.9$  & $62.7$  & $4.8$   & $74.2$  \\
\pmishift{} (ours, package) & $\mathbf{93.9}$ & $\mathbf{81.5}$ & $\mathbf{18.5}^{\dagger}$ & $\mathbf{85.7}$ \\
\midrule
$\Delta$ (\pmishift{} $-$ base) & $+4.0$ & $+18.8$ & $+13.8$ & $+11.5$ \\
$p$ (bootstrap)             & $<.001$ & $<.001$ & $<.001$ & $<.001$ \\
\bottomrule
\end{tabular}
\end{center}
\end{table}

\paragraph{Highlighted finding: the meta package \emph{mitigates}, and
does not cause, 16k degeneration.} A natural worry is that the gain is a
length artifact---that meta blocks simply buy more tokens. The budget
sweep in Table~\ref{tab:t1budget} rules this out and reveals the opposite.
\emph{Both} arms are flat from 4k to 16k (AIME: \pmishift{} $19.2\to18.5$,
base $5.0\to4.8$): extra token budget yields no accuracy on either arm, so
the treatment gap is not bought by longer generations. What \emph{does}
differ is non-termination. The base arm truncates far more at every
budget---on MATH500 $\approx$19\% of base responses hit the cap vs.\
$\approx$6\% for \pmishift{}, and on AIME $\approx$73\% vs.\
$\approx$50\%. The 16k non-termination degeneration is therefore a
property of the base policy and data that the meta mechanism
\emph{mitigates}, not one it introduces.

\begin{table}[t]
\caption{\textbf{T1b: token-budget sweep and truncation.} \emph{Preliminary,
single seed.} Accuracy is flat 4k$\to$16k on both arms (no length effect),
while the base arm truncates far more at the cap---the meta mechanism
mitigates the 16k degeneration rather than causing it.}
\label{tab:t1budget}
\begin{center}
\small
\begin{tabular}{lcccc}
\toprule
 & \multicolumn{2}{c}{AIME acc.} & \multicolumn{2}{c}{Truncation rate} \\
\cmidrule(lr){2-3}\cmidrule(lr){4-5}
Arm & 4k & 16k & MATH500 & AIME \\
\midrule
Matched base (vanilla GRPO) & $5.0$  & $4.8$  & $\approx$19\% & $\approx$73\% \\
\pmishift{} (ours, package) & $19.2$ & $18.5$ & $\approx$6\%  & $\approx$50\% \\
\bottomrule
\end{tabular}
\end{center}
\end{table}

\subsection{RQ2: Where does the effect come from? (\textsc{pending})}
\label{sec:exp:rq2}

\paragraph{Question and design.} Because T1 is a five-head package effect,
RQ2 decomposes it into pre-registered arms and adjudicates each with a
fixed decision rule. The gain could come from (i) SFT priming of meta
behaviors, (ii) the \pmishift{} reward selectively reinforcing them, (iii)
a non-semantic ``pause''/format artifact of emitting extra tokens, or (iv)
a boilerplate/shortcut shortcut of the kind AntiSD~\citep{shen2026antisd}
shows self-distillation is prone to. We instantiate four arms:
\begin{itemize}[leftmargin=1.4em,itemsep=1pt,topsep=1pt]
\item \textbf{A1 (Gandhi arm)}: meta-SFT priming $+$ \emph{vanilla} GRPO
  (correctness-only), directly instantiating
  \citet{gandhi2025cognitive}---priming without the \pmishift{} reward.
\item \textbf{A2 (shift-only arm)}: the \pmishift{} reward added on top of
  A1, but \emph{without} the calibration/format/emission maintenance
  heads---isolating the shaping reward itself.
\item \textbf{C1 (placebo, shuffled content)}: A2 evaluated with meta-block
  \emph{content} shuffled across problems, positions and lengths preserved.
\item \textbf{C2 (dedup)}: A2 with the token-set Jaccard filter enabled, so
  meta blocks that near-duplicate the preceding reasoning receive no
  credit---the AntiSD boilerplate/shortcut control.
\end{itemize}

\paragraph{Pre-registered decision rules.} These are fixed \emph{before}
the decomposition run and stated as \todo{} pending its results:
\begin{itemize}[leftmargin=1.4em,itemsep=1pt,topsep=1pt]
\item \textbf{Shift adds beyond priming} if $\text{A2} > \text{A1}$: the
  \pmishift{} reward contributes over and above meta-SFT priming amplified
  by generic RL. If $\text{A2}\approx\text{A1}$, the mechanism reduces to
  priming and the reward is inert.
\item \textbf{Content, not format} if $\text{A2} \gtrsim \text{C1}$
  (shuffled): the \emph{semantics} of the meta block, not its mere
  presence/position/length, carry the effect. If $\text{A2}\approx
  \text{C1}$, the block is a computational pause, not functional
  metacognition.
\item \textbf{Survives the AntiSD/boilerplate critique} if $\text{A2}
  \approx \text{C2}$ (dedup): the reward is not farming shortcut/boilerplate
  tokens that echo the preceding reasoning; a large drop $\text{A2}\gg
  \text{C2}$ would concede the shortcut-bias diagnosis of
  \citet{shen2026antisd}.
\end{itemize}

\begin{table}[t]
\caption{\textbf{T2 (\textsc{pending}): mechanism decomposition.} A1 (Gandhi:
meta-SFT $+$ vanilla GRPO) isolates priming; A2 (shift-only) adds the
\pmishift{} reward without maintenance heads; C1 (placebo) shuffles meta
content; C2 (dedup) enables the Jaccard boilerplate filter. Flip analysis
reports SAVE (decoy$\to$gold) and DERAIL (gold$\to$decoy) rates around meta
blocks. All cells pending; decision rules in text.}
\label{tab:t2}
\begin{center}
\small
\begin{tabular}{lcccc}
\toprule
Arm & Accuracy & SAVE rate & DERAIL rate & Net flips \\
\midrule
Matched base (no meta blocks)           & \todo{} & --      & --      & --      \\
A1 Gandhi (meta-SFT + vanilla GRPO)     & \todo{} & \todo{} & \todo{} & \todo{} \\
A2 shift-only (\pmishift{} reward)      & \todo{} & \todo{} & \todo{} & \todo{} \\
C1 placebo (shuffled content)           & \todo{} & \todo{} & \todo{} & \todo{} \\
C2 dedup (Jaccard filter on)            & \todo{} & \todo{} & \todo{} & \todo{} \\
\bottomrule
\end{tabular}
\end{center}
\end{table}

\subsection{RQ3: Difficulty stratification and OOD}
\label{sec:exp:rq3}

\paragraph{Question and design.} Metacognition should pay where the model
is actually uncertain. Table~\ref{tab:t3} stratifies held-out accuracy by
difficulty quartile (Q1 easy to Q4 hard, quartiles by the matched base's
per-problem pass rate) and reports the treatment arm's meta-emission rate
per quartile.

\paragraph{Result.} The gap is flat-to-zero where the base already
saturates and enormous where it fails. On Q1 and Q2 both arms are at
$100\%$ (no headroom). On the mid-hard Q3, \pmishift{} scores $95.7$ vs.\
the base's $84.6$ ($+11.1$\,pp), and on the hardest Q4, $46.9$ vs.\ $12.1$
($+34.8$\,pp)---nearly a $4\times$ relative improvement. The effect is thus
entirely difficulty-concentrated. Crucially, meta \emph{emission} does
\emph{not} track this: the treatment arm emits a meta block on 86--94\% of
responses across \emph{all four} quartiles (non-selective). Emission and
payoff dissociate---the model verbalizes metacognition roughly uniformly,
but that metacognition only converts to accuracy on hard problems. The
uniform emission on easy quartiles, where it cannot help, is a real
efficiency cost and a limitation (Section~\ref{sec:discussion}); the AIME
and Q4 rows double as an out-of-distribution probe, where self-distilled
metacognition transfers rather than degrading as imitation
typically does~\citep{kim2026selfdistill}.

\begin{table}[t]
\caption{\textbf{T3: difficulty-quartile stratified accuracy and
meta-emission.} \emph{Preliminary, single seed.} Quartiles by matched-base
per-problem pass rate. Q1/Q2 are saturated (no headroom); the gap grows to
$+34.8$\,pp on the hardest quartile. Emission is \emph{non-selective}
(86--94\% across all quartiles), so payoff---not emission---is
difficulty-dependent.}
\label{tab:t3}
\begin{center}
\begin{tabular}{lcccc}
\toprule
 & Q1 (easy) & Q2 & Q3 & Q4 (hard) \\
\midrule
Meta-emission rate (\pmishift{}) & \multicolumn{4}{c}{$86$--$94\%$ (non-selective across quartiles)} \\
Matched base acc.                & $100$ & $100$ & $84.6$ & $12.1$ \\
\pmishift{} acc.                 & $100$ & $100$ & $95.7$ & $46.9$ \\
$\Delta_{\mathrm{acc}}$ (\pmishift{} $-$ base)  & $0$ & $0$ & $+11.1$ & $\mathbf{+34.8}$ \\
\bottomrule
\end{tabular}
\end{center}
\end{table}

\subsection{RQ4: Calibration as a secondary outcome}
\label{sec:exp:rq4}

\paragraph{Question and design.} Our position
(Section~\ref{sec:related:calibration}) is that calibration should be
subordinate to functional metacognition, not its target. RQ4 reports
calibration for the \pmishift{} arm only, because the matched base emits
\emph{no} confidence token: with nothing to be calibrated, the base arm is
\emph{structurally uncomparable} on ECE/overconfidence. We state this
asymmetry honestly rather than fabricating a base calibration number; it is
a consequence of the protocol, not an omission. Table~\ref{tab:t4}
therefore reports per-domain and pooled 15-bin
ECE~\citep{guo2017calibration} and the overconfidence rate for the
treatment arm.

\paragraph{Result.} On the easier domains the \pmishift{} arm is
well-calibrated: ECE(15) is $0.082$ on GSM8K and $0.052$ on MATH500, with a
pooled ECE of $0.053$. On AIME, however, calibration breaks down: ECE(15)
is $0.402$ and the model is overconfident on $81.8\%$ of problems---it
states high confidence on problems it gets wrong. Accuracy (RQ1) remains
the primary outcome; this AIME miscalibration is a candid negative for the
calibration story and motivates the planned $w_{\mathrm{cal}}=0$ ablation
(Table~\ref{tab:t4}), which tests whether the \emph{functional} meta reward,
rather than the explicit Brier term, carries any calibration benefit.

\begin{table}[t]
\caption{\textbf{T4: calibration (\pmishift{} arm only).} \emph{Preliminary,
single seed.} 15-bin ECE and overconfidence rate on the held-out suite. The
matched base emits no confidence token and is \emph{structurally
uncomparable} (N/A). Calibration is a secondary outcome; accuracy
(Table~\ref{tab:t1}) is primary. AIME is badly miscalibrated (ECE $0.402$,
81.8\% overconfident).}
\label{tab:t4}
\begin{center}
\begin{tabular}{lccc}
\toprule
Arm / domain & ECE (15-bin) & Overconf.\ rate & Brier \\
\midrule
\pmishift{} --- GSM8K            & $0.082$ & \todo{} & \todo{} \\
\pmishift{} --- MATH500          & $0.052$ & \todo{} & \todo{} \\
\pmishift{} --- AIME             & $0.402$ & $81.8\%$ & \todo{} \\
\pmishift{} --- pooled           & $\mathbf{0.053}$ & \todo{} & \todo{} \\
\midrule
Matched base (no confidence)     & \multicolumn{3}{c}{N/A --- emits no confidence token (structural)} \\
\pmishift{}, $w_{\mathrm{cal}}{=}0$ (planned ablation) & \todo{} & \todo{} & \todo{} \\
\bottomrule
\end{tabular}
\end{center}
\end{table}
